\documentclass[a4paper, serif, 10pt]{moderncv}

%% moderncv
% color and style
\moderncvcolor{black}
\moderncvstyle{banking}

%% page layout/margins
\usepackage[top=2.5cm, bottom=2.5cm, left=1.5cm, right=1.5cm]{geometry}

\nopagenumbers{}

%% Babel config for German language
% ref:
% - https://latex3.github.io/babel/guides/locale-german.html
\usepackage[ngerman]{babel}

%% fontspec
\usepackage{fontspec}

\IfFontExistsTF{Linux Libertine}{
  \setmainfont[Ligatures={TeX, Common}, Numbers=OldStyle]{Linux Libertine}
}{}
\IfFontExistsTF{Linux Libertine O}{
  \setmainfont[Ligatures={TeX, Common}, Numbers=OldStyle]{Linux Libertine O}
}{}

%% CTeX
% CJK support, used to show CJK characters in the resume
%
% - heading=false: disable ctex heading
% - punct=kaiming: use kaiming punctuations styles for CJK
% - scheme=plain: use plain scheme, do not override \normalsize font size
% - space=auto: space settings for CJK characters
%
% CJK fonts are resolved per locale language, e.g. Noto Serif/Sans CJK JP
% for Japanese resumes, with fallback to the next region if not installed
%
% ref:
% - http://ctan.mirrorcatalogs.com/language/chinese/ctex/ctex.pdf
\usepackage[UTF8, heading=false, punct=kaiming, scheme=plain, space=auto]{ctex}

\IfFontExistsTF{Noto Serif CJK SC}{
  \setCJKmainfont{Noto Serif CJK SC}
}{}
\IfFontExistsTF{Noto Sans CJK SC}{
  \setCJKsansfont{Noto Sans CJK SC}
}{}

%% line spacing
\usepackage{setspace}
\setstretch{1.125}

%% URL styling - use normal text instead of monospace
\urlstyle{same}

%% Patch \httplink and \httpslink to handle full URLs
% The original moderncv macros always prepend http:// or https:// to URLs.
% This causes issues when the URL already has a protocol (e.g., https://example.com)
% resulting in malformed URLs like https://https://example.com.
% This patch checks if the URL already contains "://" and uses it directly if so.
% Note: We use \str_set:Nx (x-expansion) to expand macros like \@homepage first.
\ExplSyntaxOn
\str_new:N \g_yamlresume_url_str

\RenewDocumentCommand{\httpslink}{O{}m}{
  \str_set:Nx \g_yamlresume_url_str {#2}
  \str_if_in:NnTF \g_yamlresume_url_str {://}
    { \url{#2} }
    { \url{https://#2} }
}

\RenewDocumentCommand{\httplink}{O{}m}{
  \str_set:Nx \g_yamlresume_url_str {#2}
  \str_if_in:NnTF \g_yamlresume_url_str {://}
    { \url{#2} }
    { \url{http://#2} }
}
\ExplSyntaxOff

\name{Maximilian Brandt}{}
\title{Senior Business Analyst mit Fokus auf Prozessoptimierung, Datenanalyse und Anforderungsmanagement}
\phone[mobile]{+49 30 12345678}
\email{maximilian.brandt@example.de}
\homepage{https://www.maximilian-brandt.de}

\address{Friedrichstraße 123, Berlin, Berlin, Deutschland, 10115}{}{}

\extrainfo{{\small \faLinkedin}\ \href{https://www.linkedin.com/in/maximilian-brandt-ba}{@maximilian-brandt-ba} {} {} {} • {} {} {} 
{\small \faGithub}\ \href{https://github.com/mbrandt-analytics}{@mbrandt-analytics}}

\begin{document}

\maketitle

\section{Basisinformationen}

\cvline{}{Senior Business Analyst mit über fünf Jahren Erfahrung in der Analyse, Gestaltung und Optimierung von Geschäftsprozessen. Schwerpunkt auf der Übersetzung fachlicher Anforderungen in umsetzbare IT-Lösungen sowie auf datengetriebener Entscheidungsunterstützung.
%
\begin{itemize}
\item Erhebung, Priorisierung und Dokumentation von Anforderungen mit Stakeholdern aus Fach- und IT-Bereichen
\item Modellierung von Ist- und Soll-Prozessen mit BPMN und UML
\item Entwicklung von Kennzahlensystemen und Dashboards mit SQL, Power BI und Tableau
\item Moderation von Workshops, agilem Requirements Engineering und professioneller Stakeholder-Kommunikation
\end{itemize}}
\section{Ausbildung}

\cventry{Okt. 2018–Sept. 2020}
        {Master, Wirtschaftsinformatik, Punkte: 1,7}
        {Humboldt-Universität zu Berlin}
        {\url{https://www.hu-berlin.de}}
        {}
        {\begin{itemize}
\item Fokus auf der Verbindung von betriebswirtschaftlichen Fragestellungen und IT-Lösungen
\item Praxisprojekte zu Process Mining, Reporting und Anforderungsanalyse
\item Masterarbeit zu datengetriebener Prozessoptimierung im Finanzdienstleistungssektor
\end{itemize}
\textbf{Kurse}: Geschäftsprozessmodellierung, Datenbanken und Datenmanagement, Business Intelligence, Anforderungsmanagement, Projektmanagement, Wirtschaftsrecht}

\cventry{Okt. 2014–Sept. 2018}
        {Bachelor, Betriebswirtschaftslehre, Punkte: 1,9}
        {Technische Universität München}
        {\url{https://www.tum.de}}
        {}
        {\begin{itemize}
\item Grundlagen in Betriebswirtschaft, Statistik und quantitativer Analyse
\item Erste Projekte zu Kennzahlenanalyse und Prozessdokumentation
\end{itemize}
\textbf{Kurse}: Statistik, Finanzwirtschaft, Marketing, Operations Management, Wirtschaftsinformatik, Controlling}
\section{Berufserfahrung}

\cventry{März 2022–Heute}
        {Senior Business Analyst}
        {Nordstern Consulting GmbH}
        {\url{https://www.nordstern-consulting.de}}
        {}
        {\begin{itemize}
\item Leitung der Anforderungsanalyse für digitale Transformationsprojekte in den Bereichen Finanzdienstleistungen und E-Commerce
\item Übersetzung komplexer Fachanforderungen in User Stories, Prozessmodelle und Akzeptanzkriterien
\item Moderation von Workshops mit Stakeholdern aus Fachbereichen, IT und Management
\item Steuerung von Backlog-Priorisierung und fachlicher Qualitätssicherung in agilen Projektteams
\item Mentoring von zwei Junior Business Analysts und Etablierung einheitlicher Dokumentationsstandards
\end{itemize}
\textbf{Schlüsselwörter}: Anforderungsmanagement, Prozessoptimierung, Stakeholder-Management, Agile Methoden, Datenanalyse}

\cventry{Juli 2020–Feb. 2022}
        {Business Analyst}
        {FinTech Solutions AG}
        {\url{https://www.fintech-solutions.de}}
        {}
        {\begin{itemize}
\item Analyse und Dokumentation von Geschäftsprozessen im Online-Kreditgeschäft
\item Erstellung von Lasten- und Pflichtenheften sowie Pflege von Anforderungsspezifikationen
\item Entwicklung von SQL-Auswertungen und Power-BI-Reports für das Produktmanagement
\item Unterstützung bei Testfällen, Abnahmetests und der Einführung neuer Softwarefunktionen
\item Zusammenarbeit mit Compliance und Revision zur Sicherstellung regulatorischer Anforderungen
\end{itemize}
\textbf{Schlüsselwörter}: Anforderungsanalyse, SQL, Power BI, Testmanagement, Scrum}

\cventry{Okt. 2018–Juni 2020}
        {Junior Business Analyst}
        {Handelshaus Weber GmbH}
        {\url{https://www.handelshaus-weber.de}}
        {}
        {\begin{itemize}
\item Erstellung von Prozessdokumentationen und Auswertungen für Einkauf und Logistik
\item Pflege von Excel-basierten Reports und Vorbereitung von Management-Präsentationen
\item Unterstützung bei der Einführung eines neuen Warenwirtschaftssystems
\item Aufnahme von Anwenderfeedback und Ableitung von Verbesserungsmaßnahmen
\end{itemize}
\textbf{Schlüsselwörter}: Excel, BPMN, Prozessdokumentation, Reporting}
\section{Sprachen}

\cvline{Deutsch}{Muttersprache oder zweisprachig \hfill \textbf{Schlüsselwörter}: Muttersprache, Verhandlungssicher}
\cvline{Englisch}{Verhandlungssichere Kenntnisse \hfill \textbf{Schlüsselwörter}: C1, Business English}
\cvline{Französisch}{Gute Kenntnisse \hfill \textbf{Schlüsselwörter}: B1, Grundkenntnisse}
\section{Fähigkeiten}

\cvline{Business Analyse}{Experte \hfill \textbf{Schlüsselwörter}: Anforderungsmanagement, Prozessmodellierung, BPMN, UML, User Stories}
\cvline{Datenanalyse \& Reporting}{Erfahren \hfill \textbf{Schlüsselwörter}: SQL, Power BI, Tableau, Excel, Python (Pandas)}
\cvline{Projektmanagement \& Agile}{Erfahren \hfill \textbf{Schlüsselwörter}: Scrum, Kanban, Jira, Confluence, Stakeholder-Kommunikation}
\cvline{Domänenwissen}{Erfahren \hfill \textbf{Schlüsselwörter}: Finanzdienstleistungen, E-Commerce, Regulatorik, SAP ERP}
\section{Auszeichnungen}

\cventry{Juli 2020}
        {Humboldt-Universität zu Berlin}
        {Preis für das beste Praxisprojekt}
        {}
        {}
        {Ausgezeichnet für die Entwicklung eines datengetriebenen Dashboard-Konzepts zur Steuerung von Vertriebsprozessen im Rahmen des Masterstudiums.}
\section{Zertifikate}

\cventry{Feb. 2023}
        {IIBA}
        {Certified Business Analysis Professional (CBAP)}
        {\url{https://www.iiba.org/certification/cbap}}
        {}
        {}

\cventry{Nov. 2021}
        {Scrum.org}
        {Professional Scrum Master I}
        {\url{https://www.scrum.org/professional-scrum-certifications/professional-scrum-master-i}}
        {}
        {}

\cventry{Apr. 2022}
        {Microsoft}
        {Microsoft Certified: Power BI Data Analyst Associate}
        {\url{https://learn.microsoft.com/en-us/certifications/power-bi-data-analyst-associate/}}
        {}
        {}
\section{Veröffentlichungen}

\cventry{März 2023}
        {Anforderungsmanagement in regulierten Märkten}
        {Business Analysis Journal}
        {\url{https://www.business-analysis-journal.de/anforderungsmanagement-regulierte-maerkte}}
        {}
        {\begin{itemize}
\item Praxisorientierter Beitrag zu Methoden der Anforderungserhebung in regulierten Finanzmärkten
\item Fokus auf Traceability, Stakeholder-Alignment und dokumentierte Entscheidungswege
\end{itemize}}
\section{Projekte}

\cventry{Jan. 2023–Aug. 2023}
        {Analyse und Neugestaltung eines CRM-Prozesses für einen Finanzdienstleister}
        {CRM-Prozessanalyse und Reporting}
        {\url{https://www.nordstern-consulting.de/projekte/crm-prozessanalyse}}
        {}
        {\begin{itemize}
\item Aufnahme und Priorisierung von Anforderungen mit Fach- und IT-Abteilungen
\item Modellierung des Soll-Prozesses in BPMN und Ableitung von User Stories
\item Entwicklung eines Power-BI-Dashboards zur Steuerung von Leads und Abschlüssen
\end{itemize}
\textbf{Schlüsselwörter}: CRM, BPMN, Power BI, Anforderungsmanagement}

\cventry{März 2022–Dez. 2022}
        {Konzeption eines Kennzahlensystems für Vertrieb und Kundenservice}
        {Datengetriebene Vertriebssteuerung}
        {\url{https://www.nordstern-consulting.de/projekte/vertriebssteuerung}}
        {}
        {\begin{itemize}
\item Definition von KPIs und Datenanforderungen gemeinsam mit dem Controlling
\item Erstellung von Auswertungen und Self-Service-Reports in Tableau
\item Schulung von Fachanwendern im Umgang mit den neuen Reporting-Standards
\end{itemize}
\textbf{Schlüsselwörter}: KPI, Tableau, Datenmodellierung, Reporting}
\section{Interessen}

\cvline{Datenvisualisierung}{Tableau, Power BI, Infografiken}
\cvline{Sport}{Laufen, Radsport, Wandern}
\cvline{Musik}{Klavier, Jazz}
\section{Ehrenamtliche Tätigkeiten}

\cventry{Jan. 2022–Heute}
        {Mentor für Nachwuchs-Business-Analysten}
        {Business Analysts Berlin e.V.}
        {\url{https://www.ba-berlin.de}}
        {}
        {\begin{itemize}
\item Begleitung von Berufseinsteigern bei der Entwicklung praktischer Analysefähigkeiten
\item Organisation von Workshops zu Anforderungsmanagement und Stakeholder-Kommunikation
\item Mitwirkung an Netzwerkveranstaltungen für die Berliner Business-Analyse-Community
\end{itemize}}
    
\end{document}